\section{Conclusions}
\label{sec:conclusions}

Two improvements of the metric trial functions of unsymmetric spectral elements were
presented, analysed and tested.
\begin{enumerate}
\item \emph{Basis of the metric Vandermonde matrix.} The metric shape functions are
  basis independent. The Lagrange polynomials through the reference GLL nodes, in
  bounding-box scaled skew coordinates, give $\cond(\matr V)=1$ on affine elements and a
  condition number equal to the Lebesgue function of the metric interpolation (up to a
  bounded factor) on distorted ones. Compared with monomials this reduces the error of the
  computed shape functions by two to four orders of magnitude on distorted elements at
  $p=16$ and restores the exact symmetry of $\matr K_e$ on affine elements. The Lebesgue
  constant, and with it everything that depends only on the space, is unaffected by
  any choice of basis.
\item \emph{Hybrid trial space.} Writing only $\Pp_p(x,y)$ in metric coordinates and the
  canonical ($L^2$-orthogonal) complement in parametric coordinates retains completeness,
  affine invariance and distortion immunity, reduces to the classical element on
  straight-edged quadrilaterals, and keeps the Lebesgue constant at or close to the GLL
  value for moderate curved-edge distortion. In the tests it removed the instabilities of
  the $\Qp_p(s,t)$ element---divergence under mesh refinement in statics, negative and
  strongly complex eigenvalues, exponential growth in time---while converging with the
  optimal order. It trades this robustness for a moderately larger error constant on
  smooth static problems.
\end{enumerate}
Open questions are a stability theory (discrete inf-sup condition) for both metric spaces,
an accurate treatment of the singular integrands of the hybrid space on folded elements,
the residual (slow) growth of the hybrid element in long-time transient analysis, and the
extension to three dimensions, triangles and nonlinear problems.
